\documentclass[11pt]{article}
\usepackage[T1]{fontenc}
\usepackage{amsmath,amssymb}
\usepackage[hidelinks]{hyperref}
\usepackage[margin=1in]{geometry}

\title{Literature Status Note: ASQ Renorming and Dual-ASQ Questions}
\author{Run \texttt{fa\_banach\_001}, lane 13}
\date{2026-06-16}

\begin{document}
\maketitle

\paragraph{Status.}
\texttt{literature\_already\_answered}. This note records two exact
open-problem signals in Abrahamsen--Langemets--Lima, \emph{Almost square
Banach spaces}, arXiv:1402.0818, that are answered by later papers.

\paragraph{Source metadata.}
Trond A. Abrahamsen, Johann Langemets, and Vegard Lima,
\emph{Almost square Banach spaces}, arXiv:1402.0818. A local copy is stored as
\texttt{source\_paper.pdf}.

\paragraph{Question A: removing separability from the $c_0$ renorming theorem.}
In Section 2, immediately after Theorem 2.9 (label \texttt{thm:c0-asq};
source lines 526--561), the source proves that a separable Banach space can be
equivalently renormed to be almost square (ASQ) if and only if it contains
$c_0$, and then asks whether the same conclusion holds for a general Banach
space.

This is explicitly answered by Julio Becerra Guerrero, Gines Lopez-Perez, and
Abraham Rueda Zoca, \emph{Some results on almost square Banach spaces},
arXiv:1512.00610. Their introduction says that the main aim is to give a
positive answer to the above question. Theorem 2.2 (label
\texttt{renormageneral}) proves that every Banach space containing an
isomorphic copy of $c_0$ admits an equivalent ASQ norm. Corollary 2.3 (label
\texttt{carac0asq}) gives the exact characterization: a Banach space admits an
equivalent ASQ norm if and only if it contains an isomorphic copy of $c_0$.
A local copy is stored as \texttt{supporting\_paper\_1512.00610.pdf}.

\paragraph{Question B: can a dual space be ASQ?}
In Section 6, Question 6.6 (source lines 1438--1442), the source asks whether
the dual $X^*$ of a Banach space $X$ can be ASQ.

This is explicitly answered by Trond A. Abrahamsen, Petr Hajek, and Stanimir
Troyanski, \emph{Almost square dual Banach spaces}, arXiv:1912.12519. Their
introduction cites Question 6.6 from arXiv:1402.0818 and says that they give a
positive answer to whether ASQ spaces can be dual. Theorem 3.4 (label
\texttt{thm:bidual-ell-infty-asq}) constructs a bidual ASQ renorming of
$\ell_\infty$. Theorem 3.7 (label
\texttt{thm:dual-c0-asq-renorming}) proves the stronger characterization that
a dual Banach space admits a dual ASQ norm if and only if it contains an
isomorphic copy of $c_0$. A local copy is stored as
\texttt{supporting\_paper\_1912.12519.pdf}.

\paragraph{Scope limitation.}
This packet is a literature-status packet only. It should not be counted as a
new solution, partial theorem, or counterexample. It covers the two ASQ
renorming/dual-ASQ signals above. The separate LASQ versus WASQ separation
question in arXiv:1402.0818 is not classified by this packet.

\begin{thebibliography}{9}

\bibitem{ALL}
T. A. Abrahamsen, J. Langemets, and V. Lima,
\emph{Almost square Banach spaces}, arXiv:1402.0818.

\bibitem{BLRZ}
J. Becerra Guerrero, G. Lopez-Perez, and A. Rueda Zoca,
\emph{Some results on almost square Banach spaces}, arXiv:1512.00610.

\bibitem{AHT}
T. A. Abrahamsen, P. Hajek, and S. Troyanski,
\emph{Almost square dual Banach spaces}, arXiv:1912.12519.

\end{thebibliography}

\end{document}
